\section{Introduction}
\label{sec:introduction}

General relativity describes a dynamical spacetime geometry, while quantum
field theory describes quantum matter and its interactions on a specified
spacetime background. The aim of this work is to investigate a structure
from which both descriptions could be recovered in appropriate regimes.
The underlying objects need not themselves be points of spacetime, and
their mathematical description need not begin with a metric or an external
time parameter.

The proposal concerns a space of admissible histories together with weights
that determine how alternatives contribute to observations. A history is a
relational structure, not initially a trajectory through an already defined
space. Admissibility specifies which relations can coexist. Weighting is a
separate part of the model, since knowing which structures are possible does
not by itself determine their relative contributions. Quantum interference
requires more than an ordinary probability distribution over alternatives.
We therefore use a finite decoherence-functional description, drawing on
the histories-based treatment of quantum theory \citep{sorkin1994quantummeasure}.
This imports a quantum probability framework. It does not explain its origin.

The physical question is whether such a model can support persistent
organisation with a useful spacetime description. A successful account
would explain why the same effective laws survive a range of changes in
microscopic details. It would also have to distinguish structures that are
merely admissible from regimes that have appreciable probability within an
appropriate decoherent family. Neither persistence nor probability is
defined by resemblance to the physics one hopes to recover.

In this formulation there is no external process that traverses the space
of histories. A calculation that enumerates a history is not itself a
physical clock. Time would have to be recovered from correlations between
internal records, in the spirit of relational approaches
\citep{page1983}. This leaves the construction of clocks and the origin of
an arrow of time as separate problems. Likewise, cosmological expansion
would concern an effective metric relative to such clocks, not the number
of objects introduced by an enumeration algorithm.

Relational structures, histories, and inference principles already occur in
causal set theory \citep{bombelli1987,rideout2000csg,surya2019living},
rewriting approaches \citep{gorard2020}, and entropic dynamics
\citep{caticha2011}. These precedents do not establish the present proposal,
and no claim of novelty for these ingredients is made here. The required
comparison with those programmes remains part of the work.

To give the framework a concrete test, we specify four histories for a
two-path process and construct their weights from fixed mixing operations,
a phase, and an auxiliary route record. The model yields explicit output
probabilities, an interference visibility controlled by record overlap,
and an error bound for replacing coherent alternatives by a classical
mixture. The route and output questions separately decohere even when
their joint refinement does not. Detector uncertainty and record entropy
also differ. These calculations make the distinction between histories,
weights, and records observable within the toy model.

The example uses standard quantum mechanics. It does not explain why
quantum mechanics holds, and its relational paths are not an emergent
spacetime. Its role is to give the framework a finite instance with
explicit assumptions and consequences, rather than leave every part as an
unspecified construction.

This is a framework paper, not a derivation of a physical theory. The
methodology separates admissibility, quantum weighting, records, and
effective geometry so that assumptions in one part cannot be presented as
results of another. Section~\ref{sec:toyresults} gives the finite-model
predictions and Section~\ref{sec:discussion} states what could falsify the
prescription and what it leaves unexplained. No relational signature,
weighting law, or continuum limit is shown to reproduce general relativity
or quantum field theory.
